\SourceLeaf{162}
\begin{center}
\fontsize{8}{9.6}\selectfont\setlength{\tabcolsep}{3pt}
\begin{tabular}{l r|r}
&\(3\tfrac{1}{2}\) Proc.&4 Proc.\\
\shortstack[l]{Verpflichtungen gegen Witwen, jetzige\\und künftige Interessenten}&150210 Thl.&130877 Thl.\\
\shortstack[l]{Fonds zur Bestreitung der Bau-, Process-\\und Rechnungskosten . .}&5886 ,, &5150 ,, \\
Unverzinslicher Fonds . . . . . . . . . . . . . . . . . . . . . . . . . . . . . .&2636 ,, &2636 ,, \\\cline{2-3}
\multicolumn{1}{r}{Totalsumme . . . . .}&158732 Thl.&138663 Thl.
\end{tabular}
\end{center}
\noindent\hspace*{\parindent}Gegenüber zu stellen ist nun der Kapitalbetrag der Einnahmequellen der Witwenkasse. Da der Betrag der Beiträge der Mitglieder bei obigem ersten Ansatze bereits in Abzug gebracht war, so bleibt hier nur das Geldvermögen und die Apothekenpacht zu veranschlagen.

Das Vermögen ist in dem letzten Rechnungsabschlusse, also für 1845 Juli 1, zu 116369 Thl. 17 Ggr. 4 Pf. ausgeworfen. Um es aber auf den 1. Oct. der oben S. [149] angeführten Rechnungsgrundlage gemäss zu reduciren, muss abgezogen werden der halbjährige auf Michaelis 1845 fällig gewordene Betrag der Witwenpensionen mit 2350 Thl.; hinzugefügt hingegen die Beiträge der Mitglieder für das Jahr 1844—1845, nemlich an vollen Beiträgen 480 Thl. und an Stückzahlungen 18 Thl. 17 Ggr. 4 Pf., ferner die halbjährige Michaelis 1845 fällig gewordene Apothekenpacht mit 500 Thl. Um nichts zu übergehen, müssten auch noch die während der drei Monate 1. Juli bis 1. Oct. eingegangenen Zinsen hinzugerechnet werden, deren Betrag mir unbekannt ist; es wird aber wohl nicht viel gefehlt sein, wenn ich dafür den vierten Theil der Zinseinnahme des letzten Jahres unter Abzug der Einziehungskosten, mit 1163 Thalern in Rechnung bringe. Hienach setze ich mit Weglassung der Bruchtheile das Vermögen der Kasse am 1. Oct. 1845 zu 116171 Thalern an.

Die Apothekenpacht bei ihrer jetzigen Höhe von 1000 Thl. repräsentirt ein Kapitalvermögen von 28571 Thl. oder von 25000 Thl. je nach der Höhe des vorausgesetzten Zinsfusses. Es ist folglich das effective Vermögen der Kasse:
\begin{center}
144742 Thl. oder 141171 Thl., jenachdem man \(3\tfrac{1}{2}\) oder 4 Proc. annimmt.
\end{center}
\noindent\hspace*{\parindent}Es ergibt sich hieraus schliesslich die \textit{Bilanz} der Kasse
\begin{center}
\begin{tabular}{l}
als ein \textit{Deficit} von 13990 Thl. für Zinsfuss \(3\tfrac{1}{2}\) Proc.\\
als ein \textit{Überschuss} von 2508 Thl. für Zinsfuss 4 Proc.
\end{tabular}
\end{center}
\noindent\hspace*{\parindent}Dass der in der zweiten Voraussetzung sich ergebende Überschuss bei weitem kleiner ist, als die bei den einzelnen Bestandtheilen der Rechnung übrig bleibenden Unsicherheiten, braucht wohl kaum bemerkt zu werden: ein einziger Todesfall kann leicht auf einmal die ganze Bilanz um 3000—3500 Rthl. zum Nachtheil der Kasse abändern. Ausserdem darf man aber auch nicht übersehen, dass ich gar nichts für mögliche Verluste, durch Insolvenz der Schuldner, angesetzt habe: ich habe dieses unterlassen, weil jede präsumtive Veranschlagung höchst precär bleiben muss. Es würde interessant, aber bei der Form der Rechnungsführung der Jahresrechnungen, zumal in den frühern Zeiten, nicht ohne einen sehr grossen Zeitaufwand ausführbar sein, alle seit der Stiftung eingetretenen Verluste zusammenzustellen: ich selbst habe, um doch einigermassen eine Idee davon zu erhalten, mich mit einer Zusammenstellung der letzten 14 Jahre begnügen müssen, die (falls ich nichts übersehen habe) den Totalverlust 1729 Thl. 14 Ggr. 1 Pf., also den durchschnittlichen jährlichen Verlust \(=123\) Thl. 13 Ggr. ergeben hat. Kapitalisirt beträgt diess je nach dem Zinsfuss 3530 Thl. oder 3088 Thl., und dürfte man diesen Durchschnitt als maassgebend betrachten, so würde die \textit{Bilanz} sich für \textit{beide} Zinsfusse wie ein \textit{Deficit} herausstellen, nemlich
\begin{center}
\begin{tabular}{l}
von 17520 Thl. bei dem Zinsfuss von \(3\tfrac{1}{2}\) Procent\\
von\phantom{ 17} 580 Thl. bei dem Zinsfuss von 4 Procent.
\end{tabular}
\end{center}
\noindent Göttingen 1845 November 2.\hfill C. F. \textsc{Gauss}.
